%!TEX root = ../nauchka.tex

\section{Выбор метода параллелизма}

\begin{frame}
\frametitle{Сравнение подходов для верификации}

\begin{table}[h]
\centering
\scriptsize
\begin{tabular}{|l|c|c|c|}
\hline
\textbf{Характеристика} & \textbf{TP} & \textbf{PP} & \textbf{FSDP} \\
\hline
Шардирует веса & Да & Да & Да \\
\hline
Шардирует $A$-матрицы & Да & Нет & Нет \\
\hline
Точность bounds & Падает с глубиной & --- & Побитово идентична \\
\hline
Коммуникация & AllReduce/слой & Point-to-point & AllGather/слой \\
\hline
Совместимость с $\beta$-CROWN & Сложно & Невозможно & Да \\
\hline
\end{tabular}
\end{table}

\vspace{0.3cm}

\textbf{Вывод:} TP и FSDP --- два полюса компромисса \textit{память--точность}

\end{frame}

\begin{frame}
\frametitle{Tensor Parallelism: Column и Row Parallel}

\textbf{Column Parallel} --- веса разделены по \textbf{выходной} размерности:
\begin{itemize}
    \item $W = [W_0 | W_1]$ разделен между GPU
    \item Обратный CROWN: $A_{curr}$ уже разделена $\Rightarrow$ локально $P_i = A_i \cdot W_i$
    \item \textbf{AllReduce:} $A_{prev} = P_0 + P_1$
\end{itemize}

\vspace{0.3cm}

\textbf{Row Parallel} --- веса разделены по \textbf{входной} размерности:
\begin{itemize}
    \item Локально: $A_{prev,i} = A_{curr} \cdot W_i$
    \item Результат автоматически разделен, \textbf{коммуникация не нужна}
\end{itemize}

\vspace{0.3cm}

\begin{exampleblock}{Пара Col-Row = шардированная зона}
Между Col и Row активации имеют размерность $h/P$ на каждом GPU
\end{exampleblock}

\end{frame}

\begin{frame}
\frametitle{Проблема TP: потеря точности bounds}

\begin{alertblock}{Фундаментальное ограничение}
Внутри шардированных зон CROWN backward невозможен без полных весов $\Rightarrow$ используется IBP (Interval Bound Propagation)
\end{alertblock}

\vspace{0.3cm}

\textbf{IBP vs CROWN:}
\begin{itemize}
    \item \textbf{CROWN} --- глобальная линейная зависимость $y = Ax + b$, учитывает корреляции
    \item \textbf{IBP} --- послойные интервалы, теряет корреляции (\textit{wrapping effect})
    \item Ширина IBP-bounds растёт экспоненциально с числом слоёв
\end{itemize}

\vspace{0.3cm}

\textbf{Следствие:} TP эффективен для 1--2 шардированных зон; для глубоких моделей bounds становятся бесполезными

\end{frame}

\begin{frame}
\frametitle{FSDP для верификации}

\begin{block}{Идея}
Шардировать \textit{только веса}, собирая полную матрицу \texttt{AllGather} перед каждым слоем и освобождая после
\end{block}

\vspace{0.3cm}

\textbf{Преимущества:}
\begin{itemize}
    \item Bounds \textbf{побитово идентичны} single-GPU (каждый слой получает полную $W$)
    \item Не изменяет вычислительный граф $\Rightarrow$ совместим с $\beta$-CROWN
    \item Baseline-память (веса) сокращается на 80--90\%
\end{itemize}

\vspace{0.3cm}

\textbf{Компромисс:}
\begin{itemize}
    \item Пиковая память снижается на 34--39\% (A-матрицы полноразмерные)
    \item TP даёт $\sim\!2\times$ снижение пиковой памяти, но теряет точность
\end{itemize}

\end{frame}

\begin{frame}
\frametitle{Почему не Pipeline Parallelism?}

\begin{alertblock}{PP несовместим с верификацией}
\begin{enumerate}
    \item CROWN backward требует весов \textit{всех} слоёв одновременно
    \item При PP каждый GPU хранит только свою стадию $\Rightarrow$ пересылка весов сводит PP к менее эффективному варианту TP
    \item НЕ шардирует матрицы $A$ --- проблема памяти остаётся
\end{enumerate}
\end{alertblock}

\end{frame}
